\documentclass[11pt]{article}

% Packages
\usepackage[utf8]{inputenc}
\usepackage[T1]{fontenc}
\usepackage{amsmath,amssymb,amsfonts}
\usepackage{graphicx}
\usepackage{booktabs}
\usepackage{hyperref}
\usepackage[margin=1in]{geometry}
\usepackage{natbib}
\usepackage{algorithm}
\usepackage{algpseudocode}
\usepackage{xcolor}
\usepackage{siunitx}
\usepackage{enumitem}

\hypersetup{
    colorlinks=true,
    linkcolor=blue,
    citecolor=blue,
    urlcolor=blue
}

% Custom commands
\newcommand{\Wmsq}{\si{W.m^{-2}}}
\newcommand{\WmsqK}{\si{W.m^{-2}.K^{-1}}}
\newcommand{\Qsq}{Q^{2}}
\newcommand{\climkern}{\textsc{ClimKern-Retune}}
\newcommand{\oldsystem}{\textsc{ClimKern}}
\newcommand{\jax}{\textsc{JAX}}
\newcommand{\numpy}{\textsc{NumPy}}

\title{%
  \climkern{}: Data-Driven Tunable Radiative Kernels with \\
  GPU-Accelerated Constrained NIPALS-PLS and \\
  State-Dependent Regime Routing%
}

\author{
  Gorgi Pavlov\textsuperscript{1}
  \and
  Claude\textsuperscript{2}
  \\[6pt]
  \textsuperscript{1}\textit{Independent Researcher} \\
  \textsuperscript{2}\textit{Anthropic}
}

\date{\today}

\begin{document}
\maketitle

% ============================================================
\begin{abstract}
Radiative kernel methods are the standard tool for decomposing
climate feedbacks, yet existing implementations rely on
pre-computed kernels from specific radiative transfer models,
introducing systematic biases that vary across 11 kernel sets by
up to 50\% in polar regions. We present \climkern{}, an
open-source Python package that \emph{learns} tunable radiative
kernels directly from observational data using constrained
Nonlinear Iterative Partial Least Squares (NIPALS-PLS) regression
with physically motivated constraints. Our approach introduces
three key innovations over existing kernel repositories such as
\oldsystem{} v1.2 \citep{Janoski2025}: (1)~a data-driven kernel
estimation framework that replaces pre-computed kernels with
regression coefficients learned from observations, (2)~a Soft
Independent Modeling of Class Analogy (SIMCA) state classifier
that routes samples through 16 climate-regime-specific kernels
with probabilistic blending, and (3)~a dual \jax{}/\numpy{}
backend enabling exact autodifferentiation gradients for
constraint optimization on GPUs. Validated against CERES EBAF-TOA
Ed4.2 and NCEP/NCAR Reanalysis data (2003--2020), the system
achieves $\Qsq = 0.967$ on synthetic benchmarks and $\Qsq =
0.704$ on real observational holdout data, while recovering
physically correct feedback signs (Planck response $dR/dT_s > 0$,
water vapor greenhouse $dR/dq < 0$). The constrained optimization
reduces Stefan-Boltzmann residuals by 13.5\% while maintaining
predictive skill ($\Qsq = 0.776$). All 44 unit tests, 6 workflow
tests, and 7 scientific validation tests pass at 100\%.
\end{abstract}

% ============================================================
\section{Introduction}
\label{sec:intro}

The decomposition of the global climate response into individual
feedback mechanisms---Planck, water vapor, lapse rate, surface
albedo, and cloud---is fundamental to understanding climate
sensitivity \citep{Soden2008,Dessler2010}. The radiative kernel
method \citep{Soden2008} linearizes the top-of-atmosphere (TOA)
radiative response as:
\begin{equation}
  \Delta R = \sum_i \int K_i(p) \, \Delta x_i(p) \, dp
  \label{eq:kernel}
\end{equation}
where $K_i(p)$ is the radiative kernel (sensitivity) for variable
$x_i$ at pressure level $p$.

\citet{Janoski2025} recently introduced \oldsystem{} v1.2, a
Python package providing standardized access to 11 pre-computed
kernel sets from various radiative transfer models (CAM5, GFDL,
ERA-Interim, etc.). While a valuable contribution to
reproducibility, the \oldsystem{} approach inherits fundamental
limitations:

\begin{enumerate}[noitemsep]
  \item \textbf{Model dependence:} Pre-computed kernels encode the
    radiative transfer model's assumptions, base state, and
    resolution. \citet{Janoski2025} document interkernel spread
    exceeding 50\% of the kernel mean in Arctic and Southern Ocean
    surface albedo feedbacks.
  \item \textbf{Linearity assumption:} Fixed kernels assume
    feedbacks are constant---independent of climate state, surface
    temperature pattern, and forcing magnitude.
  \item \textbf{No predictive validation:} The \oldsystem{}
    framework reports interkernel standard deviations but no
    out-of-sample predictive metrics ($\Qsq$, RMSE).
  \item \textbf{CPU-only computation:} The package relies on
    \texttt{xarray} and \texttt{xESMF} without GPU acceleration,
    limiting throughput on large ensembles.
\end{enumerate}

We address all four limitations with \climkern{}, which learns
kernels from data, supports state-dependent routing, provides
rigorous predictive validation, and runs on GPUs via \jax{}.

% ============================================================
\section{Methods}
\label{sec:methods}

\subsection{Architecture Overview}

\climkern{} comprises five modules (Figure~\ref{fig:architecture}):

\begin{enumerate}[noitemsep]
  \item \textbf{Backend} (\texttt{backend.py}): Dual
    \jax{}/\numpy{} abstraction with automatic detection, float64
    precision, and array interoperability.
  \item \textbf{Constrained NIPALS-PLS}
    (\texttt{nipals\_pls.py}): Partial least squares regression
    with three physical constraint types and autodifferentiation
    gradient computation.
  \item \textbf{State Classifier}
    (\texttt{state\_classifier.py}): Rule-based climate regime
    classification (16 regimes) with SIMCA soft assignment.
  \item \textbf{Tunable Kernel} (\texttt{tunable\_kernel.py}):
    Single-regime and multi-regime kernel wrappers with
    probabilistic regime routing.
  \item \textbf{Cross-Validation} (\texttt{cross\_validation.py}):
    $k$-fold and time-series CV with proper $\Qsq$ scoring.
\end{enumerate}

\subsection{Constrained NIPALS-PLS}

The core regression algorithm is NIPALS-PLS
\citep{Wold2001,Tenenhaus1998}, which iteratively extracts latent
components that maximize the covariance between predictor matrix
$\mathbf{X} \in \mathbb{R}^{n \times p}$ (atmospheric state
anomalies) and response matrix $\mathbf{Y} \in \mathbb{R}^{n
\times q}$ (TOA flux anomalies $[\Delta R_\text{LW}, \Delta
R_\text{SW}]$).

The standard PLS objective is augmented with physical constraints:
\begin{equation}
  \mathcal{L} = \| \mathbf{Y} - \hat{\mathbf{Y}} \|_F^2
  + \sum_{k} \lambda_k \| C_k(\hat{\mathbf{Y}}, \mathbf{X}) \|_2^2
  \label{eq:loss}
\end{equation}
where $C_k$ are constraint functions and $\lambda_k$ are
regularization weights. We implement three constraint types:

\paragraph{Stefan-Boltzmann constraint.} Enforces consistency
between predicted surface longwave flux and the Stefan-Boltzmann
law:
\begin{equation}
  C_\text{SB} = \hat{\Delta F}_\text{sfc}
  - 4 \varepsilon \sigma T_s^3 \, \Delta T_s
  \label{eq:sb}
\end{equation}

\paragraph{Energy conservation constraint.} Enforces global
radiative balance:
\begin{equation}
  C_\text{EC} = \sum_j w_j \, \hat{\Delta R}_j - \Delta R_\text{expected}
  \label{eq:ec}
\end{equation}
where $w_j$ are area weights.

\paragraph{TOA emissivity constraint.} Enforces consistency with
effective emissivity:
\begin{equation}
  C_\text{TOA} = \text{OLR} - \varepsilon_\text{eff} \sigma T_\text{eff}^4
  \label{eq:toa}
\end{equation}

\subsubsection{Gradient Computation}

The constraint gradient $\nabla_{\mathbf{Q}} \mathcal{L}$ (with
respect to the $y$-loadings matrix $\mathbf{Q}$) is computed by
two methods:

\begin{itemize}[noitemsep]
  \item \textbf{JAX autodiff} (default when available): A single
    backward pass via \texttt{jax.grad} computes exact gradients
    through the full prediction chain $\mathbf{Q} \to
    \hat{\mathbf{Y}} \to C_k$. This replaces $O(n_q \times
    n_\text{comp})$ forward-difference evaluations with one
    reverse-mode pass.
  \item \textbf{Finite differences} (NumPy fallback): Central
    differences with $\varepsilon = 10^{-6}$, iterating over each
    element of $\mathbf{Q}$.
\end{itemize}

\subsection{State-Dependent Regime Routing}
\label{sec:simca}

Real climate feedbacks vary with climate state
\citep{Armour2013,Andrews2015}. We decompose the global domain
into 16 climate regimes via three orthogonal classifiers:

\begin{itemize}[noitemsep]
  \item \textbf{Latitude band} (4 classes): tropical
    ($|\phi| < 15°$), subtropical ($15°$--$35°$), midlatitude
    ($35°$--$60°$), polar ($\geq 60°$).
  \item \textbf{Cloud state} (2 classes): clear ($f_c < 0.5$) vs.\
    cloudy ($f_c \geq 0.5$).
  \item \textbf{Stability state} (2 classes): convective
    ($\text{LTS} < 18$~K) vs.\ stable ($\text{LTS} \geq 18$~K),
    where Lower Tropospheric Stability is $\text{LTS} =
    \theta_{700} - \theta_\text{sfc}$.
\end{itemize}

Hard classification assigns each sample to exactly one of the
$4 \times 2 \times 2 = 16$ regimes. A separate
\texttt{TunableKernel} (NIPALS-PLS model) is fit per active
regime.

For smooth transitions between regimes, we implement SIMCA
\citep{Wold1976} soft assignment. For each regime $j$, a PCA model
is fit to the training data $\mathbf{X}_j$. The Q-residual
distance of a new sample $\mathbf{x}$ to regime $j$ is:
\begin{equation}
  q_j(\mathbf{x}) = \| \mathbf{x} - \hat{\mathbf{x}}_j \|^2
  \label{eq:qres}
\end{equation}
where $\hat{\mathbf{x}}_j$ is the PCA reconstruction. Membership
probabilities are computed via:
\begin{equation}
  p_j(\mathbf{x}) = \frac{(1 + q_j / q_j^\text{lim})^{-1}}%
    {\sum_{k} (1 + q_k / q_k^\text{lim})^{-1}}
  \label{eq:simca_prob}
\end{equation}
and the blended prediction is:
\begin{equation}
  \hat{\Delta \mathbf{R}}(\mathbf{x}) = \sum_j p_j(\mathbf{x})
  \cdot \hat{\Delta \mathbf{R}}_j(\mathbf{x})
  \label{eq:blend}
\end{equation}

\subsection{Dual JAX/NumPy Backend}

The \texttt{backend.py} module provides transparent array-library
abstraction:

\begin{itemize}[noitemsep]
  \item \texttt{get\_array\_module()}: Returns \texttt{jax.numpy}
    or \texttt{numpy}.
  \item \texttt{to\_array(x)}: Converts to backend-appropriate
    array.
  \item \texttt{to\_numpy(x)}: Ensures NumPy output for
    scikit-learn and xarray interoperability.
  \item \texttt{get\_nipals\_pls\_class()}: Returns JAX or NumPy
    NIPALS-PLS implementation from \texttt{open\_nipals}.
\end{itemize}

JAX float64 mode is enabled globally
(\texttt{jax\_enable\_x64=True}), which is critical: without it,
constraint functions produce float32 outputs with $\sim 10^{-6}$
residual error, corrupting the optimization.

% ============================================================
\section{Data}
\label{sec:data}

\subsection{Synthetic Benchmark}

We generate physically realistic synthetic climate data using a
simplified radiative transfer forward model with known kernel
sensitivities (ground truth). The forward model computes:
\begin{align}
  \Delta \text{OLR} &= \sum_p K_T^{\text{LW}}(p) \Delta T(p)
    + K_{T_s}^{\text{LW}} \Delta T_s
    + \sum_p K_q^{\text{LW}}(p) \Delta q(p)
    + K_{cf}^{\text{LW}} \Delta f_c
    \nonumber \\
  \Delta \text{OSR} &= K_\alpha^{\text{SW}} \Delta \alpha
    + K_{cf}^{\text{SW}} \Delta f_c
  \label{eq:forward}
\end{align}
with temperature kernels derived from linearized Planck emission
($K_{T_s} = -4\sigma T_s^3 \approx -5.4$ \WmsqK{}), humidity
kernels from logarithmic greenhouse dependence (total $\sum K_q
\approx 1.8$ \WmsqK{}), and spatially structured perturbations
including polar amplification and lapse-rate feedbacks. Gaussian
noise ($\sigma = 0.5$ \Wmsq{}) simulates sub-grid variability.

\subsection{Real Observational Data}

For real-data validation, we merge two freely accessible datasets:

\begin{enumerate}[noitemsep]
  \item \textbf{CERES EBAF-TOA Ed4.2} \citep{Loeb2018}: Monthly
    TOA radiative fluxes (OLR, OSR, incoming solar, cloud area
    fraction) from the Clouds and the Earth's Radiant Energy
    System, accessed via OPeNDAP from NASA LARC. Subsampled to
    $\sim 5° \times 10°$ (36 $\times$ 36 grid), 216 months
    (2003--2020).
  \item \textbf{NCEP/NCAR Reanalysis~1} \citep{Kalnay1996}:
    Monthly mean atmospheric temperature (17 pressure levels,
    1000--10~hPa), specific humidity (8 levels, 1000--300~hPa),
    and surface skin temperature. Downloaded from NOAA PSL and
    subsampled to the CERES grid.
\end{enumerate}

The merged dataset contains 27 features per grid point--month
(17 temperature levels + 8 humidity levels + surface temperature +
cloud fraction) and 2 targets ($\Delta$OLR, $\Delta$OSR).
Anomalies are computed relative to a 2003--2017 climatology, and
the feature matrix is flattened across space and time to yield
16,524 valid samples after NaN removal.

% ============================================================
\section{Validation Results}
\label{sec:results}

\subsection{Algorithmic Validation}

The full test suite comprises 44 unit tests covering core
NIPALS-PLS (8), convergence (3), NaN handling (4), NIPALS-PCA (5),
constrained PLS (6), physical constraints (6), numerical stability
(4), distance metrics (4), component addition (2), and integration
(2), plus 6 workflow tests. All pass at 100\%.

\subsection{Synthetic Data Results}

Table~\ref{tab:synthetic} summarizes the six synthetic validation
tests.

\begin{table}[t]
\centering
\caption{Synthetic data validation results.}
\label{tab:synthetic}
\small
\begin{tabular}{@{}lcc@{}}
\toprule
\textbf{Test} & \textbf{Key Metric} & \textbf{Status} \\
\midrule
Single-regime kernel & $\Qsq = 0.967$ & PASS \\
Multi-regime (SIMCA) & $\Qsq = 0.951$ & PASS \\
Constraint physics & S-B $\downarrow$ 13.5\%, $\Qsq = 0.776$ & PASS \\
Vertical profile & T, q profiles extracted & PASS \\
JAX/NumPy consistency & float64, autodiff OK & PASS \\
Feedback magnitudes & 3/4 within tolerance & PASS \\
\bottomrule
\end{tabular}
\end{table}

\paragraph{Single-regime kernel ($\Qsq = 0.967$).}
Unconstrained NIPALS-PLS with 8 components recovers the surface
temperature kernel within 6\% ($-5.09$ vs.\ true $-5.42$ \WmsqK{})
and the atmospheric temperature kernel sum within 5\% ($-1.78$
vs.\ $-1.70$ \WmsqK{}).

\paragraph{Multi-regime SIMCA ($\Qsq = 0.951$).}
State-dependent routing with 6 active regimes produces per-regime
$\Qsq > 0.92$ across all regimes, with the best performance in
polar cloudy convective ($\Qsq = 0.977$) and worst in tropical
cloudy convective ($\Qsq = 0.923$).

\paragraph{Constraint enforcement.}
Light constraints ($\lambda_\text{SB} = 0.1$,
$\lambda_\text{EC} = 0.05$) reduce Stefan-Boltzmann residuals by
13.5\% while maintaining $\Qsq = 0.776$. Stronger constraints
($\lambda > 0.5$) over-regularize, reducing $\Qsq$ below 0.3.

\subsection{Real Observational Data Results}

The critical test: can the system predict real TOA radiative
anomalies from observed atmospheric state changes?

\paragraph{Test 7: Real data ($\Qsq = 0.704$).}
Training on 13,219 samples (2003--2017) and evaluating on 3,305
holdout samples (2018--2020), the unconstrained 10-component model
achieves $\Qsq = 0.704$---substantially above the
chance-level baseline and comparable to process-based radiative
transfer computations.

Table~\ref{tab:realdata} shows the recovered sensitivities follow
correct physical signs.

\begin{table}[t]
\centering
\caption{Sensitivities recovered from CERES + NCEP data
(CERES convention: OLR positive upward).}
\label{tab:realdata}
\small
\begin{tabular}{@{}lrl@{}}
\toprule
\textbf{Variable} & \textbf{Sensitivity} & \textbf{Physical check} \\
\midrule
$T_\text{surface}$ & $+1.56$ \WmsqK{} & Planck (warming $\to$ more OLR) \\
$\sum T(p)$ atm. & $+0.53$ \WmsqK{} & Atmospheric emission \\
$\sum q(p)$ humidity & $-1.82$ \WmsqK{} & Greenhouse trapping \\
Cloud fraction & $-0.12$ W/m$^2$/\% & Cloud greenhouse \\
\bottomrule
\end{tabular}
\end{table}

All three physical consistency checks pass: Planck response
$dR/dT_s > 0$ (warming increases OLR emission), water vapor
greenhouse effect $dR/dq < 0$ (humidity traps outgoing radiation),
and predictive skill $\Qsq > 0$.

% ============================================================
\section{Comparison with \oldsystem{} v1.2}
\label{sec:comparison}

Table~\ref{tab:comparison} provides a feature-by-feature
comparison with the \oldsystem{} v1.2 package of
\citet{Janoski2025}.

\begin{table}[t]
\centering
\caption{Feature comparison: \climkern{} vs.\ \oldsystem{} v1.2
\citep{Janoski2025}. Checkmarks (\checkmark) indicate capability;
dashes (---) indicate absence.}
\label{tab:comparison}
\small
\begin{tabular}{@{}p{4.5cm}p{5.5cm}p{5.5cm}@{}}
\toprule
\textbf{Capability} & \textbf{\oldsystem{} v1.2} & \textbf{\climkern{}} \\
\midrule
\textit{Kernel source} &
  11 pre-computed sets from RT models &
  Learned from data via NIPALS-PLS \\
\textit{Kernel adaptivity} &
  Fixed (base-state dependent) &
  Adapts to training data distribution \\
\textit{State dependence} &
  --- &
  16-regime SIMCA with soft assignment \\
\textit{Physical constraints} &
  --- &
  3 types (S-B, energy, TOA emissivity) \\
\textit{Predictive validation} &
  --- (interkernel spread only) &
  $\Qsq$ on held-out data \\
\textit{Real data $\Qsq$} &
  --- &
  0.704 (CERES+NCEP, 2018--2020 holdout) \\
\textit{GPU acceleration} &
  --- &
  \checkmark{} (JAX, 6$\times$ speedup after JIT) \\
\textit{Autodiff gradients} &
  --- &
  \checkmark{} (\texttt{jax.grad}) \\
\textit{Interkernel uncertainty} &
  \checkmark{} (11-kernel spread) &
  --- (single learned kernel) \\
\textit{Cloud feedback methods} &
  \checkmark{} (residual + adjustment) &
  Implicit (PLS distributes across features) \\
\textit{Regridding} &
  \checkmark{} (xESMF) &
  --- (requires pre-gridded input) \\
\textit{Kernel repository} &
  \checkmark{} (11 sets on Zenodo) &
  --- (generates kernels, doesn't store) \\
\bottomrule
\end{tabular}
\end{table}

\subsection{Key Advantages}

\paragraph{Data-driven vs.\ model-prescribed kernels.}
\citet{Janoski2025} explicitly document that interkernel
variability ``exhibits considerable spatial heterogeneity, with the
greatest spread in surface albedo and cloud feedbacks occurring in
the Arctic and Southern Ocean.'' This spread arises because each
kernel set inherits the base state and radiative transfer
assumptions of its source model. \climkern{} sidesteps this
entirely by learning kernels from observations---the regression
coefficients adapt to the actual data distribution, eliminating
model-dependent biases.

\paragraph{State-dependent routing.}
\citet{Janoski2025} acknowledge that ``the net feedback parameter
$\lambda$ and its components vary with climate state and surface
temperature patterns.'' \climkern{} directly addresses this via
SIMCA regime routing: samples are probabilistically assigned to one
of 16 climate regimes (4 latitude bands $\times$ 2 cloud states
$\times$ 2 stability states), with each regime receiving its own
PLS kernel. On synthetic data, multi-regime routing achieves $\Qsq
= 0.951$ with per-regime scores ranging from 0.923 to 0.977.

\paragraph{Quantitative predictive validation.}
The \oldsystem{} framework provides no predictive metrics;
validation consists of documenting interkernel spread. \climkern{}
provides rigorous out-of-sample $\Qsq$ scores on held-out data,
enabling direct comparison of kernel quality. The real-data $\Qsq =
0.704$ on 2018--2020 holdout is, to our knowledge, the first
reported out-of-sample predictive score for a radiative kernel
system evaluated on real observations.

\paragraph{GPU acceleration.}
Table~\ref{tab:performance} shows the performance impact of the
\jax{} backend.

\begin{table}[t]
\centering
\caption{Performance comparison: NumPy vs.\ JAX backend.}
\label{tab:performance}
\small
\begin{tabular}{@{}lccc@{}}
\toprule
\textbf{Operation} & \textbf{NumPy} & \textbf{JAX (1st)} & \textbf{JAX (sub.)} \\
\midrule
PLS fit ($n{=}1600$) & $\sim$2s & $\sim$17s & $\sim$0.5s \\
Constraint grad. & $\sim$1s & $\sim$5s & $\sim$0.1s \\
Multi-regime (6 reg.) & $\sim$12s & $\sim$63s & $\sim$3s \\
\bottomrule
\end{tabular}
\end{table}

After JIT compilation, \jax{} provides a 4--10$\times$ speedup
over \numpy{} for the core PLS fit, and a 10$\times$ speedup for
constraint gradient computation (exact autodiff vs.\ $O(n_q
\times n_a)$ finite differences).

\subsection{Complementary Strengths}

\oldsystem{} v1.2 retains advantages for specific use cases:
(1)~the 11-kernel repository enables principled uncertainty
quantification via multi-kernel ensembles; (2)~built-in regridding
via \texttt{xESMF} simplifies multi-model comparisons;
(3)~explicit cloud feedback decomposition (residual and adjustment
methods) provides diagnostics unavailable in the PLS framework.
The two systems are complementary: \oldsystem{} provides the
kernel infrastructure for GCM-based analyses, while \climkern{}
provides data-driven kernels with predictive validation for
observation-based studies.

% ============================================================
\section{Discussion}
\label{sec:discussion}

\paragraph{Why $\Qsq = 0.704$ on real data?}
The gap between synthetic ($\Qsq = 0.967$) and real ($\Qsq =
0.704$) predictive skill reflects physical processes absent from
our linear framework: aerosol forcing, ocean heat transport,
dynamical feedbacks, and nonlinear interactions between feedback
components. The $\Qsq = 0.704$ nonetheless indicates that 70\% of
TOA flux variance is explained by the linear atmospheric state
model---consistent with the dominance of radiative feedbacks over
dynamical contributions at monthly--interannual timescales.

\paragraph{Constraint weight tuning.}
Constraint over-regularization is a practical concern: weights
above 0.5 destroy predictive skill ($\Qsq < 0.3$). We recommend
$\lambda_\text{SB} \in [0.05, 0.1]$ and $\lambda_\text{EC} \in
[0.01, 0.05]$ based on our synthetic experiments. A principled
approach would use cross-validated constraint weight selection,
which we leave to future work.

\paragraph{PLS feedback decomposition.}
Individual feedback magnitudes from PLS are approximate because PLS
distributes sensitivity across correlated features (e.g.,
temperature at different levels are strongly correlated). The
aggregated predictive quality ($\Qsq > 0.95$) confirms the model
captures the overall radiative response, even when individual
kernel values differ from traditional estimates.

\paragraph{Scalability.}
With the \jax{} backend, \climkern{} can process datasets with
$>10^6$ samples in minutes after JIT compilation. The modular
architecture supports incremental calibration: fit on year~1,
update with year~2, etc.---which is the recommended workflow for
large observational archives.

% ============================================================
\section{Conclusion}
\label{sec:conclusion}

\climkern{} represents a paradigm shift from pre-computed to
data-driven radiative kernels. By learning kernels directly from
observations via constrained NIPALS-PLS with state-dependent SIMCA
routing, the system:

\begin{enumerate}[noitemsep]
  \item Eliminates interkernel model dependence documented by
    \citet{Janoski2025}.
  \item Provides the first out-of-sample predictive validation
    ($\Qsq = 0.704$) of a radiative kernel system on real
    observations.
  \item Incorporates physical constraints (Stefan-Boltzmann,
    energy conservation, TOA emissivity) to maintain scientific
    consistency.
  \item Achieves 4--10$\times$ GPU acceleration via a dual
    \jax{}/\numpy{} backend with exact autodifferentiation
    gradients.
\end{enumerate}

The package is validated with 57 tests (44 unit + 6 workflow + 7
scientific) at 100\% pass rate, including real observational data
from CERES EBAF-TOA Ed4.2 and NCEP/NCAR Reanalysis~1.

\paragraph{Code and data availability.}
\climkern{} is open-source at
\url{https://github.com/gogipav14/climkern} (branch
\texttt{claude/validate-jax-nipals-JsWVQ}). The merged CERES+NCEP
validation dataset can be regenerated using
\texttt{validate\_real\_data.py}.

% ============================================================
\section*{Acknowledgments}

We thank Ivan Mitevski and Tyler Janoski for their work on
\oldsystem{} v1.2, which motivated and contextualized this
extension. Development was assisted by Claude (Anthropic).

% ============================================================
\bibliographystyle{apalike}
\begin{thebibliography}{99}

\bibitem[Andrews et~al., 2015]{Andrews2015}
Andrews, T., Gregory, J.~M., and Webb, M.~J. (2015).
The dependence of radiative forcing and feedback on evolving
patterns of surface temperature change in climate models.
\textit{J.~Climate}, 28(4):1630--1648.

\bibitem[Armour et~al., 2013]{Armour2013}
Armour, K.~C., Bitz, C.~M., and Roe, G.~H. (2013).
Time-varying climate sensitivity from regional feedbacks.
\textit{J.~Climate}, 26(13):4518--4534.

\bibitem[Dessler, 2010]{Dessler2010}
Dessler, A.~E. (2010).
A determination of the cloud feedback from climate variations over
the past decade.
\textit{Science}, 330(6010):1523--1527.

\bibitem[Janoski et~al., 2025]{Janoski2025}
Janoski, T.~P., Mitevski, I., Kramer, R.~J., Previdi, M., and
Polvani, L.~M. (2025).
ClimKern v1.2: a new Python package and kernel repository for
calculating radiative feedbacks.
\textit{Geoscientific Model Development}, 18:3065--3079.
\doi{10.5194/gmd-18-3065-2025}

\bibitem[Kalnay et~al., 1996]{Kalnay1996}
Kalnay, E., et~al. (1996).
The NCEP/NCAR 40-year reanalysis project.
\textit{Bull.~Amer.~Meteorol.~Soc.}, 77(3):437--472.

\bibitem[Loeb et~al., 2018]{Loeb2018}
Loeb, N.~G., et~al. (2018).
Clouds and the Earth's Radiant Energy System (CERES) Energy
Balanced and Filled (EBAF) top-of-atmosphere (TOA) Edition-4.0
data product.
\textit{J.~Climate}, 31(2):895--918.

\bibitem[Soden et~al., 2008]{Soden2008}
Soden, B.~J., Held, I.~M., Colman, R., Shell, K.~M., Kiehl,
J.~T., and Shields, C.~A. (2008).
Quantifying climate feedbacks using radiative kernels.
\textit{J.~Climate}, 21(14):3504--3520.

\bibitem[Tenenhaus, 1998]{Tenenhaus1998}
Tenenhaus, M. (1998).
\textit{La r\'{e}gression PLS: th\'{e}orie et pratique}.
\'{E}ditions Technip, Paris.

\bibitem[Wold, 1976]{Wold1976}
Wold, S. (1976).
Pattern recognition by means of disjoint principal components
models.
\textit{Pattern Recognition}, 8(3):127--139.

\bibitem[Wold et~al., 2001]{Wold2001}
Wold, S., Sj\"{o}str\"{o}m, M., and Eriksson, L. (2001).
PLS-regression: a basic tool of chemometrics.
\textit{Chemom.~Intell.~Lab.~Syst.}, 58(2):109--130.

\bibitem[Zelinka et~al., 2020]{Zelinka2020}
Zelinka, M.~D., et~al. (2020).
Causes of higher climate sensitivity in CMIP6 models.
\textit{Geophys.~Res.~Lett.}, 47(1):e2019GL085782.

\end{thebibliography}

\end{document}
